% ECONOMICS_PROBLEM: {"id": "D47-95", "title": "Adaptive market maintenance", "jel_code": "D47", "source_id": "OP-0299"}
% !TeX program = xelatex
\documentclass[11pt,a4paper]{article}
\usepackage[margin=23mm]{geometry}
\usepackage{fontspec}
\setmainfont{Latin Modern Roman}
\usepackage{amsmath,amssymb,mathtools}
\usepackage{xcolor,enumitem,booktabs,longtable,array,xurl,hyperref}
\definecolor{muted}{HTML}{53636C}
\hypersetup{colorlinks=true,urlcolor=blue,linkcolor=blue}
\setlength{\parindent}{0pt}
\setlength{\parskip}{5pt}
\newcommand{\R}{\mathbb R}
\newcommand{\E}{\mathbb E}
\newcommand{\N}{\mathbb N}
\newcommand{\ind}{\mathbf 1}
\newcommand{\Var}{\operatorname{Var}}
\newcommand{\Cov}{\operatorname{Cov}}
\newcommand{\supp}{\operatorname{supp}}
\DeclareMathOperator*{\argmin}{arg\,min}
\DeclareMathOperator*{\argmax}{arg\,max}
\newcommand{\entrymeta}[3]{{\sffamily\small\color{muted}\textbf{\texttt{#1}}\quad|\quad #2\par #3\par}}
\newcommand{\Problemref}[1]{\texttt{#1}}
\newcommand{\JELref}[2]{\texttt{#2} (JEL \texttt{#1})}
\begin{document}
\section*{Adaptive market maintenance}
\textbf{Problem ID:} \texttt{D47-95}\quad \textbf{JEL:} \texttt{D47}\par
% BEGIN ECONOMICS STATEMENT
\label{problem:OP-0299}
{\sffamily\small\color{muted}Original subject group: Mechanism design and market design\par}
\label{definitions:problem:30007581}
\textbf{Audit classification:} Source/provenance unverified. NBER indexing verifies the genuine 2023 working-paper identity and DOI, but repeated full-page/PDF fetches failed. The online-control specification has no verified open-problem locator and is retained as editorial.
\subsection*{Definitions and precise research target}
\textbf{Model and research target.} A market operator repeatedly chooses a rule \(m_t\) from a finite set \(\mathcal M\), covering application limits, timing, signaling, or allocation priorities. Let \(\theta_t\) be exogenous market conditions; their evolution does not depend on selected rules. A participant state \(s_t\), including participation and adaptation, evolves as \(s_{t+1}=f(s_t,m_t,\theta_t)\), for a specified transition function \(f\) and common initial state \(s_1\). Strategic responses are incorporated into \(f\) and welfare \(W(m_t,s_t,\theta_t)\in[0,1]\). The operator observes outcomes under its chosen rule and incurs switching cost \(c(m_{t-1},m_t)\geq0\). For every comparator rule sequence \(m^*\), define its own counterfactual states by \(s^*_{t+1}=f(s^*_t,m^*_t,\theta_t)\), with \(s^*_1=s_1\). Set
\[
 R_T=\max_{m^*\text{ feasible}}\sum_t[W(m^*_t,s^*_t,\theta_t)-c(m^*_{t-1},m^*_t)]
 -\mathbb E\sum_t[W(m_t,s_t,\theta_t)-c(m_{t-1},m_t)].
\]
Both sequences obey the same announced switching budget and initial rule. For exogenous sequences satisfying \(\sum_{t=2}^T d(\theta_t,\theta_{t-1})\leq V\), with specified metric \(d\), determine minimax attainable regret over feasible adaptive policies, under explicit assumptions on state persistence, feedback, and strategic-response identification. Identify conditions allowing sublinear regret and failures caused by irreversible participation changes or unobserved counterfactuals. Specify the information needed to estimate welfare and transitions rather than treating observed outcomes as unbiased evaluations of alternatives.

\emph{Definitions and maintained assumptions (editorial unless a source definition is named).} State, exogenous-condition and action spaces, metric \(d\), transition \(f\) and welfare \(W\) must be fixed, with feedback specifying which components are observed. To compare policies, use the same exogenous sequence and initial state, and recompute each policy’s own state trajectory. If \(f\) is stochastic, fix a kernel and evaluate both welfare expectations under their own induced path laws; sharing realized endogenous states is invalid. Define comparator set \(\mathcal C_B=\{m^*_{1:T}:\#\{t\ge2:m_t^*\ne m_{t-1}^*\}\le B\}\), initial rule \(m_0\), switching-cost bounds, and learner policies with the same switch budget. Minimax policy regret is \(\inf_\pi\sup_{\theta_{1:T}:\sum d\le V}\{\sup_{m^*\in\mathcal C_B}J(m^*,\theta)-\mathbb E_\pi J(m,\theta)\}\). Declare whether adversaries are oblivious or adaptive and whether \(f,W\) are known. Even \(V=0\) can give linear regret if an early irreversible action enters an unobserved absorbing bad state; bounded exogenous variation alone is insufficient.
\subsection*{Variants and assumption changes}
\begin{enumerate}
\item \textbf{Contractive response (editorial).} Assume known \(f,W\), full-state observation and uniform contraction \(d_s(f(s,m,\theta),f(s',m,\theta))\le\kappa d_s(s,s')\) with \(\kappa<1\). Determine regret dependence on \((T,V,B)\).
\item \textbf{Unknown transitions (editorial).} Permit bandit welfare observations and an unknown finite-state kernel with a stated mixing/reset condition. Characterize exploration cost and failure without reset or counterfactual observability.
\end{enumerate}
\subsection*{References and source audit}
\begin{enumerate}
\item NBER primary indexing verifies working-paper31947 (2023) and DOI; page/PDF inaccessible. Exact online-regret question not verified. \url{https://www.nber.org/papers/w31947}
\end{enumerate}
\begin{enumerate}\raggedright
\item\hypertarget{definitions:source:30007581:1}{} Alvin E. Roth. 2023. \emph{Market Design and Maintenance}. Abstract (working paper31947); introduction in July17,2024 chapter revision.
\url{https://www.nber.org/papers/w31947}.
DOI: \url{https://doi.org/10.3386/w31947}.
\end{enumerate}

\subsection*{Common conventions: Source mathematical conventions}

These conventions apply to each entry unless explicitly replaced.

\textbf{Models and identification.} A primitive model \(m\) specifies all probability laws, feasible choices, information, equilibrium and measurement rules used by its target. The admissible class \(\mathcal M\) is fixed independently of outcomes. If \(P_O(m)\) is its observable probability law and \(\tau(m)\) the target, its identified set at observable law \(P\) is
\[
 \mathcal I_\tau(P)=\{\tau(m):m\in\mathcal M,\ P_O(m)=P\}.
\]
Point identification means that this set is a singleton. An empty set signals incompatibility of data and assumptions. Coordinate bounds need not describe the joint set. Bounds are sharp only if every retained target is attainable or arbitrarily approachable by compatible models. Finite bounds require explicit restrictions on unobserved quantities; unrestricted confounding, hidden wealth, extrapolation or arbitrary equilibrium selection can yield uninformative sets.

\textbf{Causal interventions.} A potential outcome \(Y_i(p)\) is the outcome under a complete policy regime \(p\), including timing, financing, information and market spillovers. The target population and units are fixed before treatment. With interference, use the complete assignment vector or a justified exposure mapping; individual no-interference notation alone cannot identify an economy-wide intervention. A difference of mean potential outcomes does not require the cross-world joint distribution. Conditioning on post-treatment migration, survival, enrollment or employment changes the estimand and needs separate justification.

\textbf{Statistical statements.} An estimator, confidence set, forecast or test specifies a sampling law, model class and loss/coverage criterion. Uniform \((1-\alpha)\)-coverage means \(\inf_{m\in\mathcal M}\Pr_m\{\tau(m)\in C_n\}\geq1-\alpha\), or an explicitly asymptotic version. The whole parameter space trivially has coverage, so informativeness requires a stated width, risk or power objective. A forecasting improvement is relative to a named benchmark, information set and evaluation population.

\textbf{Equilibrium and optimization.} An equilibrium comprises feasible plans, optimality, beliefs where needed, budget constraints and all market/resource clearing. The primitives or a stated benchmark define each correspondence \(\mathcal E(p;m)\). Nonemptiness is assumed or proved, never inferred just from compact policy sets. A selected-equilibrium target uses a declared selection rule; a robust target quantifies over all permitted equilibria. Use suprema or infima unless attainment is established. An \(\varepsilon\)-optimal policy is feasible and within \(\varepsilon\) of the finite optimum value. Report an empty feasible set separately.

\textbf{Welfare and accounting.} Normative utility, interpersonal normalization, social weights, numeraire, discounting and the use of public receipts are primitives. Behavioral utility need not coincide with the welfare criterion. Household welfare includes ownership income and rents once; adding profits again to compensating variation generally double-counts them. A monetary equivalent is defined by an explicit transfer and price convention, with monotonicity and range conditions. Average effects, mechanism contributions, transfers and welfare losses are different targets. Infinite sums require convergence and dynamic feasible plans require terminal or no-Ponzi conditions.

\textbf{Source versions.} A working-paper date, revision date and journal publication year are distinguished. A DOI for a journal version is not automatically the DOI of an earlier manuscript. Locators refer to the stated version and printed pages where specified. A metadata-only check does not verify a claim about a full-text conclusion. Every entry's source subsection narrows the provenance claim accordingly.
% END ECONOMICS STATEMENT
\end{document}
